% CustomVRAD — Algorithm Reference
% Compile: pdflatex algorithms.tex  (×2 for TOC/refs)
% Requires: amsmath, booktabs, hyperref, listings, microtype, tcolorbox (optional)
\documentclass[11pt,a4paper]{article}

% ─── Modern layout ───────────────────────────────────────────────────────────
\usepackage[margin=2.4cm,marginparwidth=2cm]{geometry}
\usepackage{lmodern}
\usepackage[T1]{fontenc}
\usepackage[utf8]{inputenc}
\usepackage{microtype}
\usepackage{setspace}
\setstretch{1.08}

% ─── Math & symbols ──────────────────────────────────────────────────────────
\usepackage{amsmath,amssymb,mathtools}
\usepackage{bm}

% ─── Figures, tables, code ───────────────────────────────────────────────────
\usepackage{graphicx}
\usepackage{booktabs}
\usepackage{tabularx}
\usepackage{multirow}
\usepackage{array}
\usepackage{xcolor}
\usepackage{listings}
\usepackage{enumitem}
\usepackage{caption}
\usepackage{subcaption}

% ─── Algorithms ──────────────────────────────────────────────────────────────
\usepackage{algorithm}
\usepackage{algpseudocode}

% ─── Boxes / callouts ────────────────────────────────────────────────────────
\usepackage[most]{tcolorbox}
\tcbuselibrary{breakable,skins}

% ─── Hyperlinks & meta ───────────────────────────────────────────────────────
\usepackage[hidelinks,colorlinks=true,linkcolor=cvLink,citecolor=cvLink,urlcolor=cvLink]{hyperref}
\usepackage{cleveref}

% ─── Colors ──────────────────────────────────────────────────────────────────
\definecolor{cvBlue}{HTML}{1B4F72}
\definecolor{cvAccent}{HTML}{C0392B}
\definecolor{cvMuted}{HTML}{5D6D7E}
\definecolor{cvBg}{HTML}{F4F6F7}
\definecolor{cvCodeBg}{HTML}{F8F9F9}
\definecolor{cvLink}{HTML}{1A5276}
\definecolor{cvRule}{HTML}{D5D8DC}

% ─── Title style ─────────────────────────────────────────────────────────────
\usepackage{titlesec}
\titleformat{\section}
  {\Large\bfseries\color{cvBlue}}{\thesection}{0.8em}{}[\vspace{-0.3em}{\color{cvRule}\titlerule}]
\titleformat{\subsection}
  {\large\bfseries\color{cvBlue}}{\thesubsection}{0.7em}{}
\titleformat{\subsubsection}
  {\normalsize\bfseries\color{cvMuted}}{\thesubsubsection}{0.6em}{}

% ─── Custom boxes ────────────────────────────────────────────────────────────
\newtcolorbox{notebox}[1][Note]{
  breakable, enhanced, colback=cvBg, colframe=cvBlue,
  fonttitle=\bfseries, title=#1, boxrule=0.6pt,
  left=6pt, right=6pt, top=4pt, bottom=4pt
}
\newtcolorbox{warnbox}[1][Caveat]{
  breakable, enhanced, colback=cvBg, colframe=cvAccent,
  fonttitle=\bfseries\color{cvAccent}, title=#1, boxrule=0.6pt,
  left=6pt, right=6pt, top=4pt, bottom=4pt
}
\newtcolorbox{eqbox}{
  breakable, enhanced, colback=cvCodeBg, colframe=cvRule,
  boxrule=0.4pt, left=8pt, right=8pt, top=4pt, bottom=4pt,
  sharp corners
}

% ─── Listings ────────────────────────────────────────────────────────────────
\lstset{
  basicstyle=\ttfamily\small,
  backgroundcolor=\color{cvCodeBg},
  frame=single, rulecolor=\color{cvRule},
  breaklines=true, columns=false,
  xleftmargin=4pt, xrightmargin=4pt,
  aboveskip=0.6em, belowskip=0.6em
}

% ─── Misc ────────────────────────────────────────────────────────────────────
\newcommand{\flag}[1]{\texttt{\bfseries #1}}
\newcommand{\vmt}[1]{\texttt{\$#1}}
\newcommand{\ent}[1]{\texttt{#1}}
\newcommand{\code}[1]{\texttt{#1}}
\newcommand{\RR}{\mathbb{R}}
\newcommand{\vect}[1]{\bm{#1}}

\title{%
  \vspace{-1.2em}
  {\color{cvBlue}\Huge\bfseries CustomVRAD}\\[0.35em]
  {\Large Technical Report: Lightmap Baking Algorithms}\\[0.5em]
  {\normalsize\color{cvMuted} Path-traced and radiosity pipelines for Source~SDK~2013 / Garry's~Mod}
}
\author{CustomVRAD contributors}
\date{\today}

\begin{document}
\maketitle
\begin{abstract}
\noindent
CustomVRAD is a 64-bit lightmap compiler derived from Valve VRAD (Source~SDK~2013)
for Garry's~Mod. It retains the classical patch-radiosity pipeline and optionally
replaces world-face (and static-prop) lighting with a DirectX~Raytracing (DXR)
path tracer. Extensions include hero-wavelength spectral transport, contact-hardening
soft shadows, textured emissive and filter materials, spatial volume overrides,
multi-scale ambient occlusion, IESNA photometric spot evaluation, directional
projected-texture modulation on \ent{light\_spot}, physically based underwater
media (Beer--Lambert / Kd and homogeneous volume path tracing), BSP-shipped
volumetric fog grids with convex brush SDFs and a screenspace shader,
GPU texture shadows for \vmt{alphatest}/\vmt{translucent} props, and lightmap
post-processes (denoising, seam stitching, edge pull).
This document states the algorithms as implemented: estimators, coordinate
conventions, parameter semantics, and the relationship between classic and
path-traced bake modes.
\end{abstract}

\tableofcontents
\newpage

%═══════════════════════════════════════════════════════════════════════════════
\section{System overview}
%═══════════════════════════════════════════════════════════════════════════════

CustomVRAD produces HDR and LDR lightmaps compatible with the Source engine.
World-face lighting uses exactly one of two pipelines:

\begin{enumerate}[leftmargin=1.6em]
  \item \textbf{Classic radiosity} --- stock VRAD \code{BuildFacelights} and
        iterative patch gather, optionally accelerated by OpenCL (\flag{-gpu}).
  \item \textbf{Path tracing} --- D3D12 DXR baker (\flag{-pathtrace} /
        \flag{-dxr}). On success it replaces world-face direct lighting and
        radiosity bounce. Static props may share the integrator under
        \flag{-StaticPropLighting}.
\end{enumerate}

Leaf ambient cubes, detail props, and engine-exported worldlights remain on
stock paths, with volume-aware CustomVRAD extensions where noted.
If DXR initialisation or bake fails, the compiler falls back to classic
radiosity for world faces (and CPU gather for props).

\begin{notebox}[GPU backends]
  \flag{-pt\_gpu} selects the DXR RayQuery luxel baker.
  \flag{-gpu} selects an independent OpenCL path for classic radiosity gather
  and batched AO occlusion. The two backends are not interchangeable.
\end{notebox}

\subsection{Data flow}

\begin{enumerate}[leftmargin=1.6em]
  \item Load BSP, materials (VMT/VTF), entities, static props.
  \item Build ray-tracing acceleration (CPU KD-tree and/or DXR TLAS/BLAS).
  \item Direct lighting (classic) \emph{or} PathLi bake (pathtrace).
  \item Radiosity bounce (classic only, if \flag{-bounce}~$N>0$).
  \item Optional AO, denoise, seam stitch, edge-pull-aware sampling.
  \item Write lightmaps / prop vertex colours into the BSP.
\end{enumerate}

\subsection{Notation}

\begin{tabularx}{\textwidth}{@{}lX@{}}
  \toprule
  Symbol & Meaning \\
  \midrule
  $\vect{x},\vect{n}$ & Sample position and shading normal \\
  $\rho$ & Diffuse albedo (linear RGB or spectral reflectance) \\
  $\beta$ & Path throughput \\
  $L_e$ & Emitted radiance \\
  $V$ & Binary (or soft) visibility \\
  $N$ & \flag{-pt\_bounces} --- number of \emph{indirect} hops \\
  $M$ & Path depth $= N+1$ (primary vertex $+$ $N$ bounces) \\
  \bottomrule
\end{tabularx}

%═══════════════════════════════════════════════════════════════════════════════
\section{Path tracing (DXR)}
\label{sec:pathtrace}
%═══════════════════════════════════════════════════════════════════════════════

The path tracer is a unidirectional Lambertian integrator with next-event
estimation (NEE) at every path vertex. It is implemented twice:

\begin{itemize}[leftmargin=1.4em]
  \item \textbf{GPU} (\flag{-pt\_gpu}, default): compute shader using DXR~1.1
        \code{RayQuery}, TLAS over world + optional prop BLASes
        (\code{pathtrace\_dxr\_gpu\_bake.inl}).
  \item \textbf{CPU} (\flag{-pt\_cpu}): threaded SSE reference
        (\code{PtPathLi} in \code{pathtrace\_dxr.cpp}). RGB-only; no spectral.
\end{itemize}

\subsection{PathLi estimator}

At luxel (or prop sample) $(\vect{x}_0,\vect{n}_0)$ the estimator is
\begin{equation}
  \hat{L}
  =
  \sum_{k=0}^{m}
  \beta_k\, L_{\mathrm{NEE}}(\vect{x}_k,\vect{n}_k)
  \;+\;
  \beta_{m'}\, L_{\mathrm{sky}}(\vect{x}_{m'}),
  \label{eq:pathli}
\end{equation}
where $m$ is the last surface vertex before termination, and the sky term is
added when a BSDF sample escapes to the sky (ambient sky is \emph{not} NEE-sampled).
Under \flag{-uw\_volume}, medium free-flight events may intervene between
surface vertices (\cref{sec:underwater}).

\paragraph{BSDF.}
Surfaces are Lambertian,
\begin{equation}
  f_r(\omega_i,\omega_o)=\frac{\rho}{\pi},
  \qquad
  p_{\mathrm{bsdf}}(\omega)=\frac{\cos\theta}{\pi}.
\end{equation}
After a non-sky hit the throughput update simplifies to
\begin{equation}
  \beta_{k+1} = \beta_k\,\rho_k
  \qquad\text{(albedo clamped channelwise to $[0,1]$).}
\end{equation}

\paragraph{Russian roulette.}
Starting at bounce index $k\ge 1$,
\begin{equation}
  q = \min\!\bigl(0.95,\;\max(\beta_r,\beta_g,\beta_b)\bigr)
  \quad\text{(spectral: mean of four $\beta_\lambda$)}.
\end{equation}
Terminate if $u>q$; otherwise $\beta\leftarrow\beta/q$.

\subsection{Bounce semantics}
\label{sec:bounces}

\begin{eqbox}
\begin{equation*}
  \boxed{\;\text{pathDepth}=N+1,\quad N=\texttt{-pt\_bounces}\;}
\end{equation*}
\end{eqbox}

$N$ counts \textbf{indirect hops after the luxel}. Vertex $k=0$ is always the
receiving sample (primary NEE $=$ direct lighting). Therefore:

\begin{center}
\begin{tabular}{@{}cll@{}}
  \toprule
  $N$ & Path depth & Effect \\
  \midrule
  $0$ & $1$ & Direct lights $+$ sky escape; \textbf{no GI} \\
  $1$ & $2$ & One bounce of reflected light \\
  $6$ & $7$ & Typical indoor quality (config sheet) \\
  \bottomrule
\end{tabular}
\end{center}

Defaults: $N=3$ when unset; $N=1$ under \flag{-fast}; clamp $[0,16]$.
Unset is distinguished from explicit zero by an internal sentinel ($-1$).
\flag{-pt\_prop\_bounces} mirrors this for vertex prop quality
(auto $=$ world $N$).

\subsection{Next-event estimation}

At every vertex the integrator evaluates (or samples) lights:

\begin{itemize}[leftmargin=1.4em]
  \item \textbf{Sun / sky disc} --- always evaluated; Veach power-heuristic MIS
        vs.\ cosine BSDF for the angular sun.
  \item \textbf{Point / spot} --- all lights when \flag{-pt\_lights}~$0$;
        otherwise power-proportional CDF sample of $K$ locals with weight
        $1/(p_i K)$. Bounce~$0$ always evaluates all locals on GPU.
  \item \textbf{\vmt{vrad\_emit} area triangles} --- power-weighted triangle
        selection $+$ uniform area sample (PBRT-style).
  \item \textbf{Legacy \code{emit\_surface}} --- NEE-only proxies (no BSDF hit
        contribution; avoids double-counting).
\end{itemize}

\paragraph{Sun MIS.}
For sun solid angle $\Omega=2\pi(1-\cos\theta_{\max})$,
\begin{equation}
  p_{\mathrm{light}}=\frac{1}{\Omega},
  \qquad
  w_{\mathrm{light}}
  =
  \frac{p_{\mathrm{light}}^{2}}
       {p_{\mathrm{light}}^{2}+p_{\mathrm{bsdf}}^{2}}.
\end{equation}

\paragraph{Point / spot attenuation.}
\begin{equation}
  A(d)=\frac{1}{c+\ell d+q d^{2}},
\end{equation}
with optional hard distance cap and quintic smootherstep fade
$s(t)=t^{3}(6t^{2}-15t+10)$.

\paragraph{Area emission contribution} (triangle sample with density $p_A$):
\begin{equation}
  L_{\mathrm{area}}
  =
  L_e\,
  \frac
  {(\vect{n}_r\cdot\omega)\,(-\vect{n}_e\cdot\omega)\,V}
  {d^{2}\,p_A}.
\end{equation}

\subsection{Soft shadows and contact hardening}
\label{sec:chss}

Local lights may use a soft disk of radius $r$ (\flag{-pt\_lightradius}, or
entity \ent{light\_volume} \code{Radius}). The sun uses an angular cap
(\flag{-softsun} / entity \code{SunSpreadAngle}).

Adaptive sample count (CHSS-style) from a probe blocker distance $t_{\mathrm{hit}}$
and light distance $d$:
\begin{equation}
  w_{\mathrm{pen}}
  =
  \frac{t_{\mathrm{hit}}\, r\, s_p}
       {\max(d-t_{\mathrm{hit}},10^{-3})},
  \qquad
  n
  =
  \mathrm{clamp}\!\left(
    1+\bigl\lfloor 0.35\,w_{\mathrm{pen}}+0.5\bigr\rfloor,\;
    1,\; n_{\max}
  \right),
\end{equation}
where $s_p=\texttt{-pt\_lightpenumbra}$ (default~$1$) and
$n_{\max}=\texttt{-pt\_softsamples}$ (default~$16$, clamp $1$--$64$).
Unoccluded lights use a cheap angular estimate, capped at four rays.

\flag{-pt\_softmode}~\code{direct} (default) enables soft sampling only at the
luxel; \code{all} enables it at every path vertex.

\subsection{Firefly clamp}

Hue-preserving max-channel clamp with cap $C_{\max}=2500$:
\begin{equation}
  C'
  =
  \begin{cases}
    C\cdot C_{\max}/\max(C) & \text{if }\max(C)>C_{\max},\\
    C & \text{otherwise.}
  \end{cases}
\end{equation}

\subsection{Sampling and luxel AA}

\begin{itemize}[leftmargin=1.4em]
  \item \flag{-pt\_samples}~$S$: spp per luxel (quality-derived if unset;
        clamp $1$--$4096$). GPU uses fixed $S$; CPU may adaptively raise spp
        (up to $4\times$, cap~$64$) when coefficient of variation is high.
  \item \flag{-pt\_aa}~$A$: deterministic $A\times A$ footprint grid over the
        luxel ($A\in[1,5]$, default~$3$). Spatial AA, not a multiplicative spp loop.
\end{itemize}

\subsection{Post-bake lightmap cleanup}

After PathLi accumulation the baker:
\begin{enumerate}[leftmargin=1.6em]
  \item Fills chart holes / invalid samples.
  \item Dilates suspicious dark edge outliers (sample-in-solid).
  \item Optionally denoises (\cref{sec:denoise}).
  \item Always applies two FXAA-like $3\times3$ high-contrast evening passes
        (blend starts near $25\%$ relative luminance jump, full near $55\%$,
        neighbour blend capped at $65\%$).
\end{enumerate}

%═══════════════════════════════════════════════════════════════════════════════
\section{Spectral transport}
\label{sec:spectral}
%═══════════════════════════════════════════════════════════════════════════════

Default on GPU: \flag{-pt\_spectral}. Disable with \flag{-pt\_nospectral}.
CPU PathLi remains linear RGB.

\subsection{Hero wavelengths}

Four stratified wavelengths per path over $[380,780]\,\mathrm{nm}$:
\begin{equation}
  \lambda_i
  =
  380 + (i+\xi_i)\,\frac{400}{4},
  \qquad i=0,1,2,3.
\end{equation}

\subsection{RGB $\leftrightarrow$ spectrum}

\begin{itemize}[leftmargin=1.4em]
  \item \textbf{Upsampling:} Smits~(1999) bases as tabulated in PBRT-v3
        ($32$ samples over the visible range) --- separate bases for reflectance
        and illuminants.
  \item \textbf{Transport:} illumination adds and reflectance multiplies in
        $\lambda$-space.
  \item \textbf{Projection:} Wyman analytic CIE~1931 $2^\circ$ CMFs $\to$ XYZ,
        then IEC~61966-2-1 / Rec.~709 matrix to linear sRGB.
\end{itemize}

Monte Carlo reconstruction weight uses the CIE $Y$ integral
$Y_{\mathrm{int}}=106.856895$:
\begin{equation}
  (X,Y,Z)
  \;+\!=\;
  \frac{L(\lambda)}{p_\lambda\, Y_{\mathrm{int}}}
  \bigl(\bar{x}(\lambda),\bar{y}(\lambda),\bar{z}(\lambda)\bigr).
\end{equation}

\begin{warnbox}
  Spectral mode is implemented in the GPU HLSL baker only.
  \flag{-pt\_cpu} ignores \flag{-pt\_spectral} and transports in linear RGB.
\end{warnbox}

%═══════════════════════════════════════════════════════════════════════════════
\section{Classic radiosity}
\label{sec:radiosity}
%═══════════════════════════════════════════════════════════════════════════════

Used when pathtrace is off or fails. \flag{-bounce}~$0$ skips the entire bounce
stage (direct only) --- the classic analogue of \flag{-pt\_bounces}~$0$.

\subsection{Patches and form factors}

Faces are subdivided into patches (unless bounce count is zero). Visibility is
PVS-pruned and ray-tested. Far-field differential form factor:
\begin{equation}
  F_{1\leftrightarrow 2}
  =
  \frac{(-\hat{d}\cdot\vect{n}_1)\,(\hat{d}\cdot\vect{n}_2)}
       {\|\vect{x}_1-\vect{x}_2\|^{2}}.
\end{equation}
Polygon-to-differential form factors are used when patch size invalidates the
far-field approximation. Raw transfer $=$ source area $\times$ form factor.
Receiver lists are normalised:
\begin{equation}
  s
  =
  \begin{cases}
    1/T & T>\pi,\\
    1/\pi & T\le\pi,
  \end{cases}
  \qquad
  T=\sum_j t_j.
\end{equation}

\subsection{Gather iteration}

Emitted light $E_j$ and patch reflectance $\rho_j$ gather as
\begin{equation}
  A_i = \sum_j E_j \odot \rho_j\, t_{ij}.
\end{equation}
Parent patches are area-weighted averages of children. Iteration stops after
\flag{-bounce} rounds (default~$100$) or when added energy falls below
$\max(1,\;0.0015\,E_{\mathrm{bounce\,1}})$.

\subsection{Textured bounce (\flag{-texbounce})}

Samples \vmt{basetexture} at each leaf-patch origin for colour bleed.
Reflectance is clamped to $[0,0.99]$. Supports VTF~$7.5$; falls back to flat
texdata average.

\subsection{Cavity-damped energy (\flag{-energy} / \flag{-cavity})}

Opt-in. Enclosure estimate $e=\min(1,T/\pi)$ yields
\begin{equation}
  d_E = 1 - e^{2}(1-C),
  \qquad
  d = 1+(d_E-1)\,\alpha,
\end{equation}
where $C\in[0.25,1]$ is the fully enclosed scale (default~$0.70$) and $\alpha$
is the global or volume \code{EnergyMode} blend. Open areas stay near stock;
closed rooms recirculate less.

\subsection{Artistic bounce controls}

\begin{itemize}[leftmargin=1.4em]
  \item \flag{-bounce\_boost}~$B$: scales \emph{integrated bounce only} after
        all iterations ($B\in[0,16]$; direct unchanged).
  \item \flag{-bounce\_chroma}~$C$: Oklab saturation of gathered bounce,
        strongest on early bounces:
        \begin{equation}
          (a,b)\leftarrow(a,b)\Bigl(1+\tfrac{C}{1+k}\Bigr),
          \qquad k=\text{0-based bounce index},
        \end{equation}
        keeping perceptual lightness $L$. Requires \flag{-texbounce}.
\end{itemize}

\begin{warnbox}
  Pathtrace ignores \flag{-bounce\_chroma} (non-physical). Bounce volumes may
  still apply boost to pathtrace albedo/throughput.
\end{warnbox}

\subsection{Radial splat and bounce weld}

Bounce is written to luxels with radial kernel radius scaled by
\flag{-bounce\_soft} (default~$1$, range $0.5$--$4$).

\flag{-bounce\_weld} (off by default) restricts coplanar neighbour-face GI to a
shared-edge band (normal $\cdot\ge 0.985$; falloff beyond $\sim 24$ patch radii).
Reduces rectangular GI smear across coplanar VBSP splits; can blotch ceilings.

\subsection{OpenCL gather (\flag{-gpu})}

Median-split BVH ($\le 8$ tris/leaf). Kernels batch visibility and non-bump
bounce gather over a CSR-like transfer graph (same $E_j\rho_j t_{ij}$ math).
Bump gathers stay on CPU. Default batch $32768$ rays; errors fall back to CPU.

%═══════════════════════════════════════════════════════════════════════════════
\section{Materials}
\label{sec:materials}
%═══════════════════════════════════════════════════════════════════════════════

\subsection{Emissive surfaces (\vmt{vrad\_emit*})}

Enabled by \vmt{vrad\_emit}~$1$ or \vmt{vrad\_emitstrength}${}>0$
(default strength~$200$ when flag-only). Colour from
\vmt{vrad\_emitmap}/\vmt{vrad\_emissivemap}, else \vmt{basetexture}; optional
grayscale \vmt{vrad\_emitmask}.

Texture RGB is linearised approximately as $x^{2.2}$, then scaled by strength
and \code{lightscale}. Pathtrace builds a fan-triangulated mesh and uses area
NEE; \flag{-pt\_emit\_samples}~$0$ evaluates all triangles (default~$64$
power samples). Classic mode aggregates to a softened area proxy plus a later
self-illumination pass.

\subsection{Colored glass filters (\vmt{vrad\_filter*})}
\label{sec:filter}

Pathtrace-only thin-sheet transmission. Ordinary \vmt{translucent} glass does
\emph{not} tint light unless \vmt{vrad\_filter}~$1$.

\paragraph{Transmittance construction.}
\begin{enumerate}[leftmargin=1.6em]
  \item Sample filter map (or basetexture); square RGB for linearisation.
  \item Thickness: $T\leftarrow T^{\mathrm{thickness}}$ (Beer-like artistic).
  \item Strength: Oklab lerp white $\to$ filter.
  \item Opacity: multiply by $(1-\mathrm{opacity})$.
\end{enumerate}
Rays may traverse up to $64$ filter hits, deduplicating brush-sheet
entrance/exit pairs within $12$ world units. Default two-sided; \vmt{vrad\_filter\_onesided}~$1$ = front only.

\paragraph{Stacked panes (Oklab).}
Energy transport stays linear; stacking is perceptual:
\begin{equation}
  L'=L_0 L_1,
  \qquad
  a'=a_0+a_1,
  \qquad
  b'=b_0+b_1.
\end{equation}

\subsection{Texture shadows (\flag{-textureshadows})}
\label{sec:texshadow}

Alpha-masked static props cast coverage-aware shadows without requiring the
MDL \code{STUDIOHDR\_FLAGS\_CAST\_TEXTURE\_SHADOWS} bit when
\flag{-textureshadows} is set. Materials with \vmt{alphatest} or
\vmt{translucent} (non-zero) contribute basetexture alpha into a GPU atlas
and per-triangle UV meta.

Under \flag{-pathtrace} / DXR RayQuery:
\begin{itemize}[leftmargin=1.4em]
  \item \vmt{alphatest} --- stochastic any-hit cutout (Enderton-style): a
        uniform $U\sim\mathcal{U}[0,1]$ rejects the hit when $U\ge\alpha$.
  \item \vmt{translucent} --- expected coverage transmittance
        $T\leftarrow T(1-\alpha)$ along VisRGB / SunVisRGB (soft dappled
        shadows). Binary occlusion queries still use stochastic testing.
\end{itemize}
Foliage triangles remain shadow casters only (skipped as GI bounce surfaces).
Classic SSE pathtrace-off continues to use stock CPU coverage accumulation.

%═══════════════════════════════════════════════════════════════════════════════
\section{Oklab perceptual ops}
\label{sec:oklab}
%═══════════════════════════════════════════════════════════════════════════════

Björn Ottosson's Oklab is used wherever artistic colour mixing must stay
perceptually coherent. Linear sRGB $\leftrightarrow$ Oklab follows the standard
LMS cube-root pipeline (coefficients in \code{oklab.h}).

\begin{center}
\begin{tabular}{@{}ll@{}}
  \toprule
  Use site & Operation \\
  \midrule
  Filter strength & Oklab lerp(white, $T$, strength) \\
  Stacked filters & $L$ multiply, $a,b$ add \\
  \flag{-bounce\_chroma} & Scale $a,b$; keep $L$ \\
  Env-vol inbound tint & Keep light $L$; take tint $a,b$ \\
  \bottomrule
\end{tabular}
\end{center}

NEE intensities and $\rho\times L$ products remain in linear RGB (or $\lambda$).

%═══════════════════════════════════════════════════════════════════════════════
\section{Volume entities}
\label{sec:volumes}
%═══════════════════════════════════════════════════════════════════════════════

All brush volumes share soft AABB membership: soft-min signed distance, quintic
smootherstep blend shell (\code{BlendDistance} / \code{BlendMode}), and
priority resolution (higher wins; ties $\to$ smaller AABB). Neighbour volumes
use a soft Voronoi split so outside halos do not tint each other.

\subsection{\ent{light\_env\_vol}}

Local sky / sun / ambient override. Volume lights are bake-only (not exported
as engine worldlights). Leaf ambient cubes accumulate volume-blended sun when
samples see the sun through sky. Shadow filters use the hard AABB via
\code{OutsideCastShadowIn} / \code{InsideCastShadowOut}.

Optional \code{BounceVolColor} / \code{BounceVolBright} recolour bounced light
inside the core to the volume's \code{\_light} hue.

\subsection{\ent{light\_absorb}}

Direct multiplier $1-w$ when \code{AbsorbDirect}; bounce multiplier
$1-\max(w_e,w_r)$ when \code{AbsorbBounce} (default on; strength default~$0.5$).

\subsection{\ent{light\_bounce\_vol}}

Local overrides for \code{BounceBoost}, \code{BounceChroma}, and
\code{EnergyMode}, soft-blended with CLI globals.

\subsection{\ent{light\_volume}}

Soft sphere point light: low-discrepancy samples inside radius $R$ (default~$16$).
Sample count scales with $R$ ($\sim\!16$ at default, cap~$40$; reduced with
\flag{-fast}). Early-out when fully lit/shadowed. Exported as a hard point light
at the centre (softness is bake-only).

\subsection{\ent{fog\_volume} (runtime volumetric fog)}
\label{sec:fogvol}

Bake-time construction of a 3D irradiance grid and packing of a screenspace
fog draw into the BSP (no addon). Let $\sigma$ be optical density per world
unit, $g$ Henyey--Greenstein anisotropy, and $\vect{c}$ fog tint.
Raymarching uses an AABB expanded by \code{BlendDistance} under Outside/Center
blend modes for the light grid / early-out; density membership uses a
\textbf{convex brush SDF} baked into atlas alpha from the entity's planes
(vertex-edited / clipped brushes, not only axis-aligned boxes).
Grid resolution is
\code{GridSpacing} clamped so that each axis has at most $64$ samples
($64^{3}$ cap). Optional FBM noise (\code{NoiseScale}, \code{NoiseCoverage})
modulates density; wind scrolls the atlas in world space.

\ent{fog\_volume\_blocker} carves fog (soft blend, also plane-aware) into atlas
alpha and is
rewritten to \code{info\_null}. Enabled volumes are rewritten to
\code{lua\_run\_on\_client} stubs that hook \code{RenderScreenspaceEffects}.
Clients that block map client Lua will not draw fog.

\subsection{\ent{water\_light\_vol}}

Optional override AABB for underwater IOP fitting (\cref{sec:underwater}):
Jerlov class, $\vect{\sigma}_a$, $\vect{\sigma}_s$, HG $g$, fog colour / end.
Priority resolves overlapping overrides.

%═══════════════════════════════════════════════════════════════════════════════
\section{Underwater media}
\label{sec:underwater}
%═══════════════════════════════════════════════════════════════════════════════

Physically motivated participating-medium bake for BSP water
(\code{leafWaterData} / \code{CONTENTS\_WATER}). Units: hammer inches;
ocean tables in $1/\mathrm{m}$ are scaled by $1/39.3700787$.

\subsection{Modes}

\begin{center}
\begin{tabular}{@{}lll@{}}
  \toprule
  Flag & Mode & Estimator \\
  \midrule
  \flag{-underwater} & A & Beer--Lambert segment $T$ $+$ Kd downwelling for sky/sun \\
  \flag{-uw\_volume} & C & Homogeneous free-flight volume PT (implies A media) \\
  \flag{-uw\_off} & --- & Force disable \\
  \flag{-uw\_maxscatter}~$N$ & C cap & Max medium events (default~$8$, max~$64$) \\
  \bottomrule
\end{tabular}
\end{center}

Both CPU PathLi and the GPU RayQuery baker implement the modes; Mode~A also
scales classic \code{GatherSampleLight} for direct sun/sky/local paths.

\subsection{IOP construction from VMT fog}

Source \vmt{fogend} is authored for \emph{screenspace fog}, not lightmap
optical depth. Fit:
\begin{equation}
  \sigma_{\mathrm{lum}}
  =
  \frac{-\ln 0.05}{\texttt{fogend}}
  \cdot s_{\mathrm{bake}},
  \qquad
  s_{\mathrm{bake}}=0.28,
\end{equation}
channel-biased by inverse fog colour so transport shifts toward the fog tint.
Scatter fraction $0.45$ of extinction; diffuse downwelling
\begin{equation}
  \vect{k}_d = 0.35\,\vect{\sigma}_t
\end{equation}
(Kd $\ll$ beam extinction). Optional \vmt{vrad\_uw\_jerlov},
\vmt{vrad\_uw\_sigma\_a}/\vmt{vrad\_uw\_sigma\_s}, \vmt{vrad\_uw\_g},
or \ent{water\_light\_vol} overrides. Default fog end when missing: $800$.

\subsection{Mode A (attenuation)}

Segment transmittance through each water AABB of length $L$:
\begin{equation}
  T=\exp(-\vect{\sigma}_t L)
  \quad\text{(componentwise)}.
\end{equation}
Sky/sun irradiance at a submerged receiver uses \emph{only}
\begin{equation}
  S_{\mathrm{sky}}(\vect{x})=\exp(-\vect{k}_d\, z_{\mathrm{depth}}),
\end{equation}
floored at $0.04$ per channel --- not $T\times S_{\mathrm{sky}}$
(double counting). Local lights use segment $T$ only. Mode~C disables Kd.

\subsection{Mode C (volume path tracing)}

Inside water, free-flight distance $t=-\ln(1-u)/\max(\sigma_t)$ is accepted
only if $t$ is strictly less than the AABB exit distance; otherwise the ray
continues as a surface path (no fake boundary scatter). On a medium event,
throughput is multiplied by $T$ to the event and by scatter albedo
$\vect{\sigma}_s/\vect{\sigma}_t$, HG direction is resampled, and a small
sky-veil term reinjects ambient fill. Escapes to sky apply segment $T$
through the water column (Kd off).

\begin{warnbox}
  Dense Source water materials with short \vmt{fogend} still darken deep
  pools; author longer fog ends, Jerlov classes, or \ent{water\_light\_vol}
  IOPs for clear water.
\end{warnbox}

%═══════════════════════════════════════════════════════════════════════════════
\section{Ambient occlusion}
\label{sec:ao}
%═══════════════════════════════════════════════════════════════════════════════

Multi-scale hemisphere AO in \code{FinalLightFace}, enabled by \flag{-ao} /
\flag{-ao\_*} or entities \ent{light\_ao} / \ent{light\_ao\_vol}.

\begin{equation}
  AO(\vect{x})
  =
  \frac{\sum_i w_i V_i}{\sum_i w_i},
  \qquad
  w_i=\max(\vect{n}\cdot\omega_i,\,0.08).
\end{equation}
Directions with $\vect{n}\cdot\omega<-0.25$ are rejected (small wrapped
hemisphere for contact darkening). Ray origin:
\begin{equation}
  \vect{x}_0
  =
  \vect{x}+\vect{n}\,b+\omega\max(0.05,\,0.5 b).
\end{equation}
Final multiplier $M=\max(0,\,1-S+S\cdot AO)$ with strength $S\in[0,8]$.

Defaults: samples~$16$, distance~$48$, strength~$1$, bias~$0.25$.
Optional bilateral luxel denoise:
\begin{equation}
  w
  =
  \exp\!\Biggl[
    -\frac{\|\Delta x\|^{2}}{2\sigma_s^{2}}
    -\frac{(AO_j-AO_i)^{2}}{2(0.12)^{2}}
  \Biggr].
\end{equation}
With \flag{-gpu}, AO can use batched OpenCL occlusion.

\begin{warnbox}[Pathtrace interaction]
  Under \flag{-pathtrace} with $N=\texttt{-pt\_bounces}>0$, post-multiply AO is
  \emph{skipped}: the path tracer already accounts for geometric occlusion.
  \flag{-ao\_force} re-enables the multiply for artistic double-darkening.
  With $N=0$ (direct$+$sky only), AO still applies as a stand-in for missing GI.
\end{warnbox}

%═══════════════════════════════════════════════════════════════════════════════
\section{Soft sun}
\label{sec:softsun}
%═══════════════════════════════════════════════════════════════════════════════

Replaces stock's fixed $30$ random cone rays:

\begin{itemize}[leftmargin=1.4em]
  \item Low-discrepancy golden-angle spiral in the sun cone.
  \item Adaptive count $\sim\!9$ at $2^\circ$ up to $\sim\!40$ at $50^\circ$.
  \item Early-out after a short probe when fully lit or fully shadowed.
\end{itemize}
$0^\circ$ $=$ hard sun. Typical soft look: $0.5^\circ$--$3^\circ$.

%═══════════════════════════════════════════════════════════════════════════════
\section{Photometric spots and projected textures}
\label{sec:photometric}
%═══════════════════════════════════════════════════════════════════════════════

Bake-only extensions on stock \ent{light\_spot}. Engine-exported worldlights
retain cone keys for dynamic lighting; lightmap evaluation may substitute an
IESNA intensity table and/or a projected Valve Texture Format (VTF) multiplier.

\subsection{IESNA photometry}
\label{sec:ies}

When the entity key \code{IES} names an LM-63 file under \texttt{garrysmod/IES/}
(or a per-map subdirectory \texttt{IES/maps/\textit{mapname}/}), the angular candela table
replaces \code{\_inner\_cone}, \code{\_cone}, and \code{\_exponent} for the bake.
Only \texttt{TILT=NONE} profiles are supported. Distance attenuation
(\code{\_constant\_attn}, \ldots) is unchanged.

Let $\vect{b}$ be the beam axis and $\vect{r}$ a right vector from entity angles.
For a unit direction $\vect{\omega}$ from light to sample, vertical and horizontal
angles $(\theta,\phi)$ index a peak-normalized table $I(\theta,\phi)\in[0,1]$.
The photometric weight is $w_{\mathrm{IES}}=s\,I(\theta,\phi)$ with scale
$s=\texttt{IESScale}$ (default~$1$). Optional \code{IESBrightness}${}\ge 0$
overrides stock \code{\_light} brightness for the bake only;
\code{IESMaxIntensity} caps per-sample RGB after attenuation.

\subsection{Projected texture classification}
\label{sec:projtex}

The key \code{ProjectedTexture} names a material under \texttt{materials/}
(optional \texttt{materials/} prefix and \texttt{.vtf} suffix are stripped).
The loader classifies the VTF:

\begin{center}
\begin{tabular}{@{}lll@{}}
  \toprule
  Kind & Detection & Domain \\
  \midrule
  Planar (2D) & No cubemap / no usable six-face layout & Spot cone \\
  Cubemap & Six faces (\code{IsCubeMap} / ENVMAP + face count) & Full sphere \\
  Spherical 2D & ENVMAP without six faces & Full sphere \\
  \bottomrule
\end{tabular}
\end{center}

Spherical 2D maps use equirectangular sampling when width $\ge\tfrac{3}{2}$ height;
otherwise an OpenGL-style spheremap / matcap projection. Cubemap VTFs (including
VTF~$7.5$ with CRC resources) are decoded to RGBA face arrays; GPU upload stores
six consecutive \code{Texture2DArray} layers per cubemap in Source
\code{CUBEMAP\_FACE\_*} order ($+X,-X,+Y,-Y,+Z,-Z$).

\subsection{Planar framing in the outer cone}

For planar projections, let $\alpha=\arccos(\texttt{\_cone})$ (outer half-angle)
and $\vect{\omega}$ the light-to-sample direction. With $\cos\theta=\vect{\omega}\cdot\vect{b}$:
\begin{equation}
  u'=\frac{\vect{\omega}\cdot\vect{r}}{\cos\theta\,\tan\alpha},
  \qquad
  v'=\frac{\vect{\omega}\cdot\vect{u}}{\cos\theta\,\tan\alpha},
\end{equation}
where $\vect{u}$ completes a right-handed basis. Texture coordinates are
\begin{equation}
  (u,v)=\bigl(\tfrac12+\tfrac12\,f\,u',\;
              \tfrac12+\tfrac12\,f\,v'\bigr),
\end{equation}
with framing factor $f=1$ for \code{ProjectedTextureMode=fill} (square sides on
the cone; corners clipped) and $f=\sqrt{2}$ for \code{fit} (square inscribed in
the cone). Samples with $(u,v)\notin[0,1]^{2}$ contribute zero (hard mask).
Cubemap and spherical modes ignore the cone and evaluate over $4\pi$.

\subsection{Cubemap direction mapping}
\label{sec:cubemap}

Entity angles define a light-local frame $(\vect{r},\vect{u},\vect{b})$
(right, up, beam). A world direction $\vect{\omega}$ maps to Source cube axes
(Z-up) by
\begin{equation}
  \vect{d}
  =
  \bigl(
    \vect{\omega}\cdot\vect{r},\;
    \vect{\omega}\cdot\vect{b},\;
    \vect{\omega}\cdot\vect{u}
  \bigr),
\end{equation}
i.e.\ $(d_x,d_y,d_z)=(\mathrm{right},\,\mathrm{forward},\,\mathrm{up})$.
Face selection and face coordinates follow the Direct3D cubemap table
(identical to hardware \code{texCUBE} / Source VTF layout):
\begin{equation}
\begin{aligned}
  +X&:\; sc=-d_z,\; tc=-d_y, &
  -X&:\; sc=+d_z,\; tc=-d_y, \\
  +Y&:\; sc=+d_x,\; tc=+d_z, &
  -Y&:\; sc=+d_x,\; tc=-d_z, \\
  +Z&:\; sc=+d_x,\; tc=-d_y, &
  -Z&:\; sc=-d_x,\; tc=-d_y,
\end{aligned}
\end{equation}
with major-axis magnitude $m$ and
\begin{equation}
  u=\tfrac12\bigl(sc/m+1\bigr),
  \qquad
  v=\tfrac12\bigl(tc/m+1\bigr).
\end{equation}
Faces clamp at edges. CPU and GPU bakers share this convention.

\subsection{Composite intensity}

Let $L_0$ be the attenuated spot intensity (photometric or cone),
$\vect{P}\in[0,1]^{3}$ the projected RGB sample, and $w_{\mathrm{IES}}$ the
optional IES weight. The bake contribution uses
\begin{equation}
  L = L_0\, w_{\mathrm{IES}}\, \vect{P}
\end{equation}
(componentwise; missing IES or projection $\Rightarrow$ factor~$1$).
When both IES and \code{ProjectedTexture} are set, IES acts solely as an
angular mask on the projection. Evaluation is wired through classic gather,
CPU PathLi, GPU PathLi, and static-prop lighting.

%═══════════════════════════════════════════════════════════════════════════════
\section{Denoising}
\label{sec:denoise}
%═══════════════════════════════════════════════════════════════════════════════

Opt-in via \flag{-pt\_denoise}; select backend with \flag{-pt\_denoiser}.

\subsection{OIDN (default)}

Intel Open Image Denoise \code{RTLightmap} model on CPU. One thread-local
filter per bake worker (pinned to avoid oversubscription). Charts smaller than
$16^{2}$ are edge-clamp padded.

\subsection{OptiX}

NVIDIA CUDA HDR denoiser (no albedo/normal guides). Shared serialised GPU
resources; same $16^{2}$ padding.

\subsection{Sakai (statistical)}

Sakai et~al.\ (2024)--inspired Welch-gated repeated $5\times5$ filter using
variance of the mean where available:
\begin{equation}
  t
  =
  \frac{|L_i-L_j|}
       {\sqrt{\mathrm{Var}(\bar L_i)+\mathrm{Var}(\bar L_j)
              +(0.04 L_{\mathrm{ref}})^{2}+\varepsilon}}.
\end{equation}
If $t>1.35$, neighbour weight is attenuated by
$\exp\!\bigl[-1.5\bigl((t-1.35)/1.35\bigr)^{2}\bigr]$.
Runs $1+\texttt{radius}$ passes, then two low-gradient $3\times3$ evening
passes. \flag{-pt\_denoise\_radius} / \flag{-pt\_denoise\_strength} apply here
(OIDN/OptiX largely ignore radius).

%═══════════════════════════════════════════════════════════════════════════════
\section{Static prop lighting}
\label{sec:props}
%═══════════════════════════════════════════════════════════════════════════════

\subsection{Classic speedups}

Four-wide SSE direct gather; with \flag{-gpu}, bounce rays that miss / hit sky
skip the BSP lightmap walk.

\subsection{Pathtrace props}

Under \flag{-pathtrace} $+$ \flag{-StaticPropLighting}:

\begin{itemize}[leftmargin=1.4em]
  \item \textbf{Prop lightmap texels} --- same spp / bounces / NEE as world.
  \item \textbf{Vertex lighting} --- with \flag{-pt\_prop\_vertgrid}~$G>0$
        (default~$4$): each LOD0 triangle gets an $8\times8$ chart packed into a
        $64\times64$-cell atlas page ($512\times512$), inset by $1$ texel
        (right-triangle lattice in the $6\times6$ content region,
        $\approx21$ samples/tri). Pages are hole-filled and, when
        \flag{-pt\_denoise} is on, denoised with the world-face backends.
        Each vertex takes the closest valid sample from every adjacent
        triangle chart and averages those colours. Unlit verts inherit from
        nearest lit. $G=0$ $=$ old one-sample-per-vertex.
\end{itemize}
Override quality with \flag{-pt\_prop\_samples} /
\flag{-pt\_prop\_bounces} (defaults: $\max(8,S/4)$ spp; bounces $=$ world).
Multi-LOD models share lit colours across LODs.

%═══════════════════════════════════════════════════════════════════════════════
\section{Lightmap sample quality}
\label{sec:luxelq}
%═══════════════════════════════════════════════════════════════════════════════

\subsection{Edge pull}

Default on (inset $0.5$ luxels; range $0$--$2$ via \flag{-edgepull}).
For chart width $W$, sample coordinate $s$ is clamped to $[i,\,W-1-i]$
(similarly $t$). Moves edge samples off adjacent solids --- thin walls, jambs,
floor/wall scallops. \flag{-noedgepull} restores stock positions.

\subsection{Seam stitching}

On by default after \code{FinalLightFace}. Coplanar faces that share an edge
(VBSP splits) blend near-edge luxels toward the neighbour value at the same
world position so both sides converge to a common average. Per lightstyle and
bump layer; displacements skipped. Disable with \flag{-nostitch}.

%═══════════════════════════════════════════════════════════════════════════════
\section{Reference integrator (pseudocode)}
\label{sec:pseudocode}
%═══════════════════════════════════════════════════════════════════════════════

\begin{algorithm}[H]
\caption{PathLi --- one spp at a surface sample}
\begin{algorithmic}[1]
\Require position $\vect{x}$, normal $\vect{n}$, bounce limit $N$, seed
\State $\beta \leftarrow \mathbf{1}$;\; $L\leftarrow\mathbf{0}$;\;
       $\mathrm{depth}\leftarrow N+1$
\For{$k\leftarrow 0$ \textbf{to} $\mathrm{depth}-1$}
  \State $L \leftarrow L + \beta\cdot\mathrm{NEE}(\vect{x},\vect{n},k)$
         \Comment{soft only if $k=0$ or softmode=all}
  \State $\omega\leftarrow\mathrm{CosineHemisphere}(\vect{n})$
  \State trace closest hit along $\omega$, applying filter panes to $\beta$
  \If{miss or sky}
    \State $L\leftarrow L+\beta\cdot L_{\mathrm{sky}}(\vect{x})$;\;
           \textbf{break}
  \EndIf
  \State $\beta\leftarrow\beta\cdot\rho_{\mathrm{hit}}$
         \Comment{or spectral Reflect($\lambda$)}
  \If{$k\ge 1$}
    \State Russian roulette on $\beta$
  \EndIf
  \State $\vect{x},\vect{n}\leftarrow$ hit point / oriented normal
\EndFor
\State \Return firefly-clamp$(L)$
\end{algorithmic}
\end{algorithm}

%═══════════════════════════════════════════════════════════════════════════════
\section{Parameter cheat sheet}
\label{sec:params}
%═══════════════════════════════════════════════════════════════════════════════

\begin{center}
\small
\begin{tabular}{@{}llp{5.8cm}@{}}
  \toprule
  Flag & Default & Role \\
  \midrule
  \flag{-pathtrace} & off & Enable DXR baker \\
  \flag{-pt\_gpu}/\flag{-pt\_cpu} & GPU & Integrator backend \\
  \flag{-pt\_samples}~$S$ & quality & spp per luxel \\
  \flag{-pt\_bounces}~$N$ & $3$ ($1$ fast) & Indirect hops; $0$=direct+sky \\
  \flag{-pt\_aa}~$A$ & $3$ & Luxel footprint grid \\
  \flag{-pt\_spectral} & on & Hero-$\lambda$ GPU transport \\
  \flag{-pt\_nospectral} & --- & Linear RGB path transport \\
  \flag{-pt\_lights}~$K$ & $0$=all & Local NEE samples (bounce$>0$) \\
  \flag{-pt\_emit\_samples} & $64$ & Area-emit NEE draws ($0$=all) \\
  \flag{-pt\_lightradius} & $0$ & Soft disk for point/spot \\
  \flag{-pt\_lightpenumbra} & $1$ & CHSS scale \\
  \flag{-pt\_softsamples} & $16$ & Soft ray cap \\
  \flag{-pt\_softmode} & direct & Soft at luxel / all vertices \\
  \flag{-pt\_denoise} & off & Enable denoise \\
  \flag{-pt\_denoiser} & oidn & oidn $\mid$ optix $\mid$ sakai \\
  \flag{-pt\_prop\_vertgrid} & $4$ & Prop vert atlas on ($0$=per-vert) \\
  \flag{-underwater} & off & Mode~A Beer--Lambert $+$ Kd \\
  \flag{-uw\_volume} & off & Mode~C free-flight volume PT \\
  \flag{-uw\_maxscatter}~$N$ & $8$ & Cap medium scatters (Mode~C) \\
  \flag{-bounce}~$N$ & $100$ & Classic radiosity rounds ($0$=direct) \\
  \flag{-texbounce} & off & Textured radiosity albedo \\
  \flag{-energy}/\flag{-cavity} & off / $0.70$ & Cavity damp \\
  \flag{-bounce\_boost} & $1$ & Artistic bounce scale \\
  \flag{-bounce\_chroma} & $0$ & Early-bounce Oklab sat \\
  \flag{-bounce\_weld} & off & Coplanar edge GI weld \\
  \flag{-gpu} & off & OpenCL classic gather \\
  \flag{-textureshadows} & off & Prop $\alpha$ shadows (\vmt{alphatest} / \vmt{translucent}) \\
  \flag{-ao} & off & Hemisphere AO pass \\
  \flag{-ao\_force} & off & AO even under pathtrace GI \\
  \flag{-edgepull} & $0.5$ & Luxel edge inset \\
  \flag{-nostitch} & --- & Disable seam stitch \\
  \bottomrule
\end{tabular}
\end{center}

\begin{notebox}[Config presets]
  Hammer-friendly flag sheets live under \texttt{configs/} and load via
  \flag{-config}~\texttt{full} (searches \texttt{.cfg}/.txt next to the binary).
  Do not put \flag{-game} / map path inside the config file.
\end{notebox}

%═══════════════════════════════════════════════════════════════════════════════
\section{Implementation map}
\label{sec:files}
%═══════════════════════════════════════════════════════════════════════════════

\begin{center}
\small
\begin{tabular}{@{}ll@{}}
  \toprule
  Module & Responsibility \\
  \midrule
  \code{pathtrace\_dxr.cpp} & CPU PathLi, bake orchestration, AA, cleanup \\
  \code{pathtrace\_dxr\_gpu\_bake.inl} & HLSL PathLi (spectral, NEE, filters) \\
  \code{pathtrace\_dxr\_device.*} & D3D12 / DXR device, AS, dispatches \\
  \code{pt\_spectral.h} & Smits bases, CIE CMFs, RGB$\leftrightarrow\lambda$ \\
  \code{oklab.h} & Perceptual colour ops \\
  \code{vrad\_emit*.cpp} & Textured emission + area mesh \\
  \code{vrad\_filter.cpp} & Colored glass transmittance \\
  \code{vrad.cpp} / \code{vismat.cpp} & Classic radiosity loop \\
  \code{radial.cpp} & Bounce splat / weld \\
  \code{vrad\_gpu.*} & OpenCL BVH + gather \\
  \code{ao.cpp} & Multi-scale AO \\
  \code{absorb.cpp} / \code{bounce\_vol.cpp} & Volume effects \\
  \code{fog\_volume.*} & Fog grid bake + BSP Lua/shader pack \\
  \code{water\_medium.*} & Underwater IOPs, Beer--Lambert, volume PT hooks \\
  \code{ies\_profile.*} & IESNA LM-63 load + angular eval \\
  \code{light\_projection.*} & Projected VTF load, CPU sample, GPU arrays \\
  \code{pathtrace\_*denoise*} & OIDN / OptiX / Sakai \\
  \code{vradstaticprops.cpp} & Prop SSE + pathtrace vert atlas \\
  \code{lightmap.cpp} & Edge pull, final face lighting, spot parse \\
  \code{loadcmdline.cpp} & \flag{-config} preset tokenization \\
  \bottomrule
\end{tabular}
\end{center}

%═══════════════════════════════════════════════════════════════════════════════
\section*{References \& provenance}
%═══════════════════════════════════════════════════════════════════════════════

\begin{itemize}[leftmargin=1.4em]
  \item Kajiya, J.~T.\ --- The rendering equation (SIGGRAPH 1986).
  \item Veach, E.\ --- Robust Monte Carlo methods for light transport (PhD, 1997):
        power-heuristic MIS.
  \item Smits, B.\ --- An RGB to spectrum conversion for reflectances (JGTools 1999);
        PBRT-v3 tabulated bases.
  \item Wyman, C.\ et~al.\ --- Analytic CIE~1931 CMF fits.
  \item Ottosson, B.\ --- Oklab (2020).
  \item Pharr, Jakob, Humphreys --- \emph{Physically Based Rendering} (area lights, spectra).
  \item Sakai et~al.\ --- Statistical lightmap filtering (2024 inspiration for the Sakai denoiser).
  \item Intel Open Image Denoise --- \code{RTLightmap} model.
  \item NVIDIA OptiX --- HDR AI denoiser.
  \item IESNA LM-63 --- standard file format for luminous intensity distributions.
  \item Microsoft Direct3D --- cubemap face selection and $(s,t)$ mapping
        (\code{texCUBE} / environment maps).
  \item Jerlov / Solonenko \& Mobley --- ocean optical water types (IOP tables).
  \item Jaffe--McGlamery --- classical underwater image formation (direct /
        forward / backscatter decomposition).
  \item Monzon et~al.\ (2024) --- spectral underwater radiative transfer
        (downwelling $K_d$).
  \item Henyey \& Greenstein --- phase function for volume scattering.
  \item Valve Source SDK 2013 --- stock VRAD radiosity / lightmap pipeline;
        VTF cubemap face order (\code{CUBEMAP\_FACE\_*}).
\end{itemize}

\vspace{1.5em}
{\color{cvMuted}\footnotesize
\noindent
Algorithms are described as implemented in the CustomVRAD source tree.
CLI defaults and clamps may change between revisions; \texttt{vrad.exe} help
output and \texttt{configs/full.cfg} define shipping presets.
}

\end{document}
